% !TeX document-id = {0860a59e-3926-43bd-9926-0464d3bf14f1}
% !TEX program = LuaLaTeX --shell-escape
\documentclass[10pt,a6paper]{article}
\pagestyle{empty}

% 1. MARGÓK BEÁLLÍTÁSA
\usepackage{geometry}
\geometry{
	top=1cm,
	bottom=1cm,
	left=1.5cm,
	right=1.5cm,
}

% betűtípus
\usepackage{fontspec}
\setmainfont{Linux Libertine O} % Linux Libertine betűtípus használata

% Nyelv beállítása (latin)

\usepackage{polyglossia}
\setdefaultlanguage{latin} % Alapértelmezett nyelv latin

% 3. Gregoriotex CSAK EZUTÁN
\usepackage{gregoriotex} 

% Csomag a hasábokhoz (Ez hiányzott!)
\usepackage{multicol}

% Címsorok számozásának kikapcsolása
\setcounter{secnumdepth}{-1}

% --- GREGORIO BEÁLLÍTÁSOK ---
% Alapméret (ha 17-es kell, itt állítsd be)
\grechangestaffsize{15}

% Iniciálé
\grechangedim{beforeinitialshift}{1.5mm}{scalable}
\grechangedim{afterinitialshift}{1.5mm}{scalable}
\grechangestyle{initial}{\fontsize{36}{36}\selectfont}

% Kottaszín és feliratok
\gresetlinecolor{black}
\gresetheadercapture{commentary}{grecommentary}{string}
\grechangestyle{annotation}{\small\bfseries}

% Térközök
\grechangedim{baselineskip}{42pt}{scalable}
\grechangedim{spaceabovelines}{5mm}{scalable} % Frissítve 5mm-re a korábbi kódod alapján
\grechangedim{spacebeneathtext}{1.5mm}{scalable} 
\grechangedim{interglyphspace}{0.9mm}{scalable}

%Sorkizárt
\gresetlastline{justified}

%Modositojelek dominikanus kezelese
\gresetalterationeffect{line}

% --- CÍMSOROK BEÁLLÍTÁSA (titlesec javítva) ---
\usepackage[compact]{titlesec}

% Térközök
\titlespacing*{\section}{0pt}{*0}{*1}
\titlespacing*{\subsection}{0pt}{*0}{*1}
\titlespacing*{\subsubsection}{0pt}{10pt}{*0}

% Formázás (KOMA-Script parancs kiváltása a section-nél)
\titleformat{\section}{\normalfont\LARGE\scshape\center}{\thesection}{1em}{}
\titleformat{\subsection}{\normalfont\Large\center}{\thesubsection}{1em}{}
\titleformat{\subsubsection}{\normalfont\normalsize\itshape}{\thesubsubsection}{1em}{}

\begin{document}
	
	% Helyesebb cím formázás az article osztályban
	\section*{Zsoltártónusok}
	
	Az \emph{Éneklő Egyház} és a \emph{Graduale Hungaricum} magyar
	sorozata alapján.
	
	\subsubsection{1. zsoltártónus} \gregorioscore[a]{scores/tonus/ee_tonus1_alap}
	\begin{multicols}{2}
		\subsubsection{1 re} \gregorioscore[a]{scores/tonus/ee_1D_term}
		\subsubsection{1 lá} \gregorioscore[a]{scores/tonus/ee_1g_term}
	\end{multicols}
	
	\subsubsection{2. zsoltártónus}
	\gregorioscore[a]{scores/tonus/ee_tonus2}
	
	\subsubsection{3. zsoltártónus} \gregorioscore[a]{scores/tonus/ee_tonus3}
	
	\subsubsection{4. zsoltártónus} \gregorioscore[a]{scores/tonus/ee_tonus4}

	\subsubsection{5. zsoltártónus}
	\gregorioscore[a]{scores/tonus/ee_tonus5}
	
	\subsubsection{6. zsoltártónus}
	\gregorioscore[a]{scores/tonus/ee_tonus6}
	
	\subsubsection{7. zsoltártónus} \gregorioscore[a]{scores/tonus/ee_tonus7_alap}
	\begin{multicols}{2}
		\subsubsection{7 lá} \gregorioscore[a]{scores/tonus/ee_7a_term}
		\subsubsection{7 dó} \gregorioscore[a]{scores/tonus/ee_7c_term}
	\end{multicols}
	
	% Szintaktikai hiba javítva (hiányzó zárójel)
	\subsubsection{8. zsoltártónus} \gregorioscore[a]{scores/tonus/ee_tonus8_alap}
	\begin{multicols}{2}
		\subsubsection{8 szó} \gregorioscore[a]{scores/tonus/ee_8G_term}
		\subsubsection{8 dó} \gregorioscore[a]{scores/tonus/ee_8c_term}
	\end{multicols}
	
	\subsubsection{tonus peregrinus}
	\gregorioscore[a]{scores/tonus/ee_tonus_peregrinus}
\end{document}
